% ======================================================================
%  Section 4.6.2 --- Controlled conditions.  Corrected. KEEP.
%
%  A. A CONTRADICTION WITH 4.6.4 THAT THE RESULTS EXPOSE.
%     This subsection lists "the stopping criterion" among the quantities
%     held fixed but does not say what it is. Section 4.6.4 then defines
%     failure as not reaching the stopping threshold within the iteration
%     cap. The threshold in the implementation is ||e||^2 > 10^4, and
%     4.7.1 reports the feature error settling three to four orders of
%     magnitude above it. No run reaches it. By 4.6.4's definition every
%     run in the chapter failed. The criterion has to be stated here, and
%     4.6.4's definition of convergence has to be the one 4.7.1 actually
%     uses.
%
%  B. P IS NEVER GIVEN A VALUE. It is introduced in 4.2.1 and used in
%     4.3.2 (the breakdown point), 4.4.3 (where N was needed alongside it)
%     and 4.7.3 (percentages of the measurements). This subsection is where
%     the configuration is declared and is the place to state it.
%
%  C. "the occluded target texture" is held fixed across the three
%     variants but not across the two test configurations, the nominal case
%     having no occluder at all. As written it reads as though one texture
%     serves throughout.
%
%  D. The initial and desired poses are listed as held fixed but never
%     given. They belong in 4.6.3, which is currently an empty heading.
% ======================================================================

\subsection{Controlled conditions}
\label{sec:ablation-conditions}

Within each test configuration, everything other than $\mathbf{D}$ is held
fixed across the three variants: the Hermite filter bank and its four scales,
the interaction matrix~\eqref{eq:stacked-L}, the adaptive gain and
damping~\eqref{eq:lambda}--\eqref{eq:mu}, the floor~\eqref{eq:hessian-floor}
on the damped Hessian, the sampling period, the stopping criterion, the
iteration cap, and the initial and desired poses. The target texture is fixed
across the variants of a given configuration and is what distinguishes the two
configurations of Section~\ref{sec:test-configurations} from each other.
%%% (C) "within each test configuration" at the front, and the texture
%%% named as the thing that varies BETWEEN configurations rather than
%%% listed among the constants.
The implementation, the simulation platform~\cite{key18} and the camera
configuration are those described in Section~\ref{sec:her-sim-config} of
Chapter~\ref{ch:herpvs}. The weighting is therefore the sole independent
variable.

%%% (B) The dimensions, stated once, where the configuration is declared.
The images are $320\times240$ pixels and the filters have $9\times9$ support,
so the ten-pixel border they cannot reach leaves an interior of
$220\times300=66\,000$ positions carrying measurements. With four scale
responses stacked at each position the feature vector has dimension
$P=264\,000$, and the quantity $N$ of~\eqref{eq:Gbar} is $66\,000$. The
percentages of measurements quoted in Section~\ref{sec:weight-maps} are
fractions of $P$; the structure measure is defined on the $N$ positions.

%%% (A) The criterion, stated, with its consequence.
The stopping criterion inherited from the implementation halts a run when the
feature error falls below an absolute threshold,
$\norm{\mathbf{e}}^{2}<10^{4}$, or when the iteration cap is reached. As
Section~\ref{sec:nominal} reports, the feature error settles three to four
orders of magnitude above that threshold in every run, so the absolute
criterion is never met and every run proceeds to the cap; results are read at
iteration~600, by which point the commanded velocity has fallen below
$10^{-3}\,\mathrm{m\,s^{-1}}$ in all cases. Convergence is therefore assessed
throughout this chapter against the commanded velocity rather than against the
inherited threshold, and Section~\ref{sec:metrics} is to be read accordingly.
%%% Without this paragraph, 4.6.4's "a run being declared to have failed if
%%% the stopping threshold is not reached within the iteration cap" declares
%%% every run in the chapter a failure.

Two points of procedure are recorded for reproducibility. The Tukey tuning
constant $c$ is held at the same value in Variants~2 and~3, so that the two
differ in the argument supplied to the weight function and not in the function
itself. And the robust scale estimate is computed, in each variant, from the
residuals that variant actually submits to the estimator: from $e_{j}$ in
Variant~2 and from $\tilde{e}_{j}$ in Variant~3.

% ======================================================================
%  E. CONSEQUENT EDIT TO 4.6.4 -- required, not optional
%
%  4.6.4 currently reads, in part:
%     "... and whether the run converged at all, a run being declared to
%      have failed if the stopping threshold is not reached within the
%      iteration cap."
%  Replace with something of the form:
%     "... and whether the run converged, a run being taken as converged
%      once the commanded velocity falls below 10^{-3} m/s, for the reason
%      given in Section 4.6.2."
%  It also records "the number of iterations to convergence" as a metric.
%  That number is only defined once convergence is defined; either
%  recompute it against the velocity criterion or drop it, since the
%  chapter does not in fact report it anywhere.
%
%  F. 4.6.3 "Test configurations" is still an empty heading and is where
%  the poses belong. It needs, at minimum:
%     nominal   : reference and servoing images both from the clean texture
%     occluded  : reference from the clean texture, servoing images from
%                 the occluded one, disc covering 7.4% of the measurements
%                 at the initial pose and 5.80% at the desired pose
%     both      : t_0 = (40, -40, 1100) mm, omega_0 = (12, -12, 6) deg
%                 t* = (0, 0, 1300) mm,     omega* = (0, 0, 0) deg
%                 initial displacement about 41 px mean, 65 px worst case
%  I can draft it if you send me the heading.
%
%  G. NOT CHANGED
%  The two procedural points at the end are the strongest part of the
%  subsection. Holding c fixed and computing each scale from the residuals
%  the variant actually submits are precisely the two things a sceptical
%  reader would check, and stating them pre-empts the check.
% ======================================================================
